% ---------------------------------------------------------------------
% CHAPTER 2 — BACKGROUND
% Job: give the reader the shared vocabulary the RQs will need.
% Kept tight; this is not a survey.
% ---------------------------------------------------------------------
\chapter{Background}

\section{Software architecture, drift, and conformance checking}
% Reflexion modelling / SACC as the practical motivation for mapping.

\section{From dependency graphs to representation learning for SAR}
% SDG assumption; clustering vs search vs fusion; then GNNs/GAE/embeddings.

\section{Graph neural networks and encoders used in this thesis}
% GCN, GAT, VGAE/GAE, HGT — only what's needed to read the contributions.
% Homophily / heterophily; message-passing reach; why encoder != the bound.

\section{Evaluation: metrics, ground truth, and their limits}
% MoJoFM/A2A/C2Ccvg/ARI/TurboMQ; precision-recall-F1 for mapping.
% Why we read metrics jointly; GT is human-made and contestable; effect sizes
